% PAGES: 285-285
% Opens subsection 9.4.1. Ends mid-derivation: the align block for (9.66)
% closes with a trailing comma, continuing on PDF page 286 (folio 270) with
% "since, from (9.60), it follows that". Footnote 7 (referenced by
% "(9.63),$^7$" below) is attached here since its marker appears on this
% page; its text is transcribed at the point where it is anchored.

\subsection{Maximum return portfolios and the tangency portfolio}

As was the case for the minimum variance portfolio, see section 9.3.1, any maximum
return portfolio consists of a cash position that is determined to match the required
variance of the portfolio return, with the rest of the portfolio invested in the
tangency portfolio.

Recall from (9.36) that the asset weights vector in the tangency portfolio can be
written as
\begin{equation}
w_T \;=\; \frac{1}{\bfone^t\Sigma_R^{-1}\bar\mu}\Sigma_R^{-1}\bar\mu,
\end{equation}
where $\frac{1}{\bfone^t\Sigma_R^{-1}\bar\mu}$ is a constant and $\Sigma_R^{-1}\bar\mu$
is a vector. Then,
\begin{align}
w_T^t\Sigma_Rw_T &= \frac{1}{\left(\bfone^t\Sigma_R^{-1}\bar\mu\right)^2}
\left(\Sigma_R^{-1}\bar\mu\right)^t\Sigma_R\left(\Sigma_R^{-1}\bar\mu\right) \notag\\
&= \frac{1}{\left(\bfone^t\Sigma_R^{-1}\bar\mu\right)^2}
\bar\mu^t\left(\Sigma_R^{-1}\right)^t\left(\Sigma_R\Sigma_R^{-1}\right)\bar\mu \notag\\
&= \frac{1}{\left(\bfone^t\Sigma_R^{-1}\bar\mu\right)^2}\bar\mu^t\Sigma_R^{-1}\bar\mu,
\end{align}
where for (9.63) we used the facts that $\left(\Sigma_R^{-1}\right)^t=\Sigma_R^{-1}$
since $\Sigma_R^{-1}$ is a symmetric matrix, see Lemma 5.4, and that
$\Sigma_R\Sigma_R^{-1}=I$.

By taking the square root on both sides of (9.63),\footnote{Note that
$\bar\mu^t\Sigma_R^{-1}\bar\mu > 0$ since $\Sigma_R$ is a symmetric positive definite
matrix, and therefore $\Sigma_R^{-1}$ is also symmetric positive definite; cf. Lemma
5.4.} we find that
\[
\sqrt{w_T^t\Sigma_Rw_T} \;=\; \frac{1}{\bfone^t\Sigma_R^{-1}\bar\mu}
\sqrt{\bar\mu^t\Sigma_R^{-1}\bar\mu}\cdot
\operatorname{sign}\!\left(\bfone^t\Sigma_R^{-1}\bar\mu\right),
\]
and therefore
\begin{equation}
\frac{\bfone^t\Sigma_R^{-1}\bar\mu}{\sqrt{\bar\mu^t\Sigma_R^{-1}\bar\mu}} \;=\;
\frac{1}{\sqrt{w_T^t\Sigma_Rw_T}}\cdot
\operatorname{sign}\!\left(\bfone^t\Sigma_R^{-1}\bar\mu\right).
\end{equation}

Then, from (9.60) and (9.64), it follows that
\begin{equation}
w_{max,cash} \;=\; 1-\frac{\sigma_P}{\sqrt{w_T^t\Sigma_Rw_T}}\cdot
\operatorname{sign}\!\left(\bfone^t\Sigma_R^{-1}\bar\mu\right).
\end{equation}

Moreover, from (9.59) and (9.62), we find that
\begin{align}
w_{max} &= \frac{\sigma_P}{\sqrt{\bar\mu^t\Sigma_R^{-1}\bar\mu}}\Sigma_R^{-1}\bar\mu
\notag\\
&= \sigma_P\frac{\bfone^t\Sigma_R^{-1}\bar\mu}{\sqrt{\bar\mu^t\Sigma_R^{-1}\bar\mu}}
\cdot\frac{1}{\bfone^t\Sigma_R^{-1}\bar\mu}\Sigma_R^{-1}\bar\mu \notag\\
&= \sigma_P\frac{\bfone^t\Sigma_R^{-1}\bar\mu}{\sqrt{\bar\mu^t\Sigma_R^{-1}\bar\mu}}\,w_T
\notag\\
&= (1-w_{max,cash})\,w_T,
\end{align}
